%%%%%%%%%%%%%%%%%%%%%%%%%%%%%%%%%%%%%%%%%%%%%%%%%%%%%%%%%%%%%%%%%%%%%%
% LaTeX Example: Project Report
%
% Source: http://www.howtotex.com
%
% Feel free to distribute this example, but please keep the referral
% to howtotex.com
% Date: March 2011 
% 
%%%%%%%%%%%%%%%%%%%%%%%%%%%%%%%%%%%%%%%%%%%%%%%%%%%%%%%%%%%%%%%%%%%%%%
% How to use writeLaTeX: 
%
% You edit the source code here on the left, and the preview on the
% right shows you the result within a few seconds.
%
% Bookmark this page and share the URL with your co-authors. They can
% edit at the same time!
%
% You can upload figures, bibliographies, custom classes and
% styles using the files menu.
%
% If you're new to LaTeX, the wikibook is a great place to start:
% http://en.wikibooks.org/wiki/LaTeX
%
%%%%%%%%%%%%%%%%%%%%%%%%%%%%%%%%%%%%%%%%%%%%%%%%%%%%%%%%%%%%%%%%%%%%%%
% Edit the title below to update the display in My Documents
%\title{Project Report}
%
%%% Preamble
\documentclass[paper=a4, fontsize=11pt]{scrartcl}
\usepackage[T1]{fontenc}
\usepackage{fourier}

\usepackage[english]{babel}															% English language/hyphenation
\usepackage[protrusion=true,expansion=true]{microtype}	
\usepackage{amsmath,amsfonts,amsthm} % Math packages
\usepackage[pdftex]{graphicx}	
\usepackage{url}
\usepackage{titlesec}
\usepackage{mathtools}
\usepackage{enumitem}
%\usepackage{lstlisting}


%%% Custom headers/footers (fancyhdr package)
\usepackage{fancyhdr}
\pagestyle{fancyplain}
\fancyhead{}											% No page header
\fancyfoot[L]{}											% Empty 
\fancyfoot[C]{}											% Empty
\fancyfoot[R]{\thepage}									% Pagenumbering
\renewcommand{\headrulewidth}{0pt}			% Remove header underlines
\renewcommand{\footrulewidth}{0pt}				% Remove footer underlines
\setlength{\headheight}{13.6pt}
\usepackage[margin=0.5in]{geometry}

%%% Equation and float numbering
\numberwithin{equation}{section}		% Equationnumbering: section.eq#
\numberwithin{figure}{section}			% Figurenumbering: section.fig#
\numberwithin{table}{section}				% Tablenumbering: section.tab#

%% Section Format
\titleformat*{\section}{\LARGE\bfseries}
\titleformat*{\subsection}{\Large\bfseries}
\titleformat*{\subsubsection}{\large\bfseries}
\titleformat*{\paragraph}{\large\bfseries}
\titleformat*{\subparagraph}{\large\bfseries}

\DeclarePairedDelimiter{\ceil}{\lceil}{\rceil}

%%% Maketitle metadata
\newcommand{\horrule}[1]{\rule{\linewidth}{#1}} 	% Horizontal rule

\title{
		%\vspace{-1in} 	
		\usefont{OT1}{bch}{b}{n}
		\normalfont \normalsize \textsc{University of Purdue\\
      Data Communication and Computer Networks\\
      CS 53600 Spring 2020} \\ [25pt]
		\horrule{0.5pt} \\[0.4cm]
		\huge Assignment 6\\
		\horrule{2pt} \\[0.5cm]
}
\author{
		\normalfont\normalsize Pedro da Costa Abreu Júnior\\ [-5pt]
    \normalsize Student ID 0030878383\\[-5pt]
    \normalsize \today
}
\date{}


%%% Begin document
\begin{document}
\maketitle

\section*{Problem 1}
\begin{enumerate}
  \item CSMA/CA, because we cannot perform collision detection due to the hidden
terminal problem. CDMA can also be a good option.
  \item CSMA/CD, because the probability of collision is fairly low, but when it
happens it is possible to recover with exponential backoff.
  \item CSMA/CA or CDMA for the same reasons as 1.
\end{enumerate}

\section*{Problem 2}
\begin{enumerate}
  \item No,
    Src IP: 192.168.3.001
    Dest IP: 192.168.3.003
    Src Mac: 77-77-77-77-77-77
    Dest Mac:99-99-99-99-99-99
  
  \item Yes,
    Src IP: 192.168.3.001
    Dest IP: 192.168.2.003
    Src Mac: 77-77-77-77-77-77
    Dest Mac:88-88-88-88-88-88,
    At least three datagrams will be needed to deliver it to B.
    First E must deliver a datagram to Router 2, then Router 2 must deliver it to Router 1 and then finally Router 1 will deliver it to B.

  \item Yes, Router 1 will receive the ARP broadcast from A. However, Router 1 will not forward this ARP request to subnet 2. B will not send a ARP request to A because the broadcast will contain A MAC address. Router 1 will not receive the data frame from B to A. Assuming the that the switch forwarding table is initially empty, it will contain A and B IP and MAC addresses.

\end{enumerate}

\section*{Problem 3}
First the host needs to get an IP address by sending a DHCP broadcast. Then, the DHCP server will respond with a DHCP ACK message, containing an allocated IP address for the host use, the IP address of the DNS server, the IP address of the default gateway router, and a network mask. The host then contacts the DNS server by sending its message to the gateway router. The DNS server responds with the IP address C, that the host wishes to download the page from. With C address, the host can finally create a TCP socket. After the TCP three-way handshake the host can send an HTTP GET to C, and C will respond with an HTTP response message, with the requested page in the body of the message. And then C sends the HTTP message through the socket. When the message is finally received by the host it can extract the Web page and finally display it to the user.


%%% End document
\end{document}

%%% Local Variables:
%%% mode: latex
%%% TeX-master: t
%%% End:
